% TOP_PROBLEM: {"id": 1107, "title": "Casas-Alvero Conjecture", "category_id": 4, "category": {"id": 4, "name": "algebra", "display_name": "Algebra", "description": "Group theory, ring theory, field theory, and algebraic structures.", "slug": "algebra", "order_index": 4, "created_at": "2026-07-31T15:26:25.671Z"}, "status": "open"}
% FIRST_POSTED: null
% ENTRY_KIND: statement_only
% Imported from: batch04.tex; UnsolvedMath ID 1107
\documentclass[11pt]{article}
\usepackage[margin=25mm]{geometry}
\usepackage{fontspec}
\setmainfont{Latin Modern Roman}
\usepackage{amsmath,amssymb,mathtools}
\usepackage{unicode-math}
\setmathfont{Latin Modern Math}
\usepackage{xurl,hyperref}
\hypersetup{colorlinks=true,urlcolor=blue}
\setcounter{secnumdepth}{0}
\setlength{\parindent}{0pt}
\setlength{\parskip}{5pt}
\setlength{\emergencystretch}{3em}
\newcommand{\sref}[2]{\hyperlink{source:#1:#2}{[S#2]}}
\newcommand{\cataloguescope}[1]{\paragraph{Recorded target.}#1}
\begin{document}
% UnsolvedMath ID 1107; source catalogue code ALG-007
\setcounter{section}{1106}
\section{Casas-Alvero Conjecture}\label{1107}
{\small\sffamily UnsolvedMath ID 1107 · Algebra}
\subsection{Definitions and mathematical statement}

\paragraph{Shared notation.}$\mathbb N=\{1,2,\ldots\}$, and $\mathbb Z,\mathbb Q,\mathbb R,\mathbb C$ denote the integers, rationals, reals and complex numbers. Any other conventions are specified locally.


Let \(K\) be any field of characteristic zero.  For \(f\in K[X]\), write
\(f^{(j)}\) for its \(j\)-th ordinary formal derivative, and let
\(\gcd(f,f^{(j)})\) be a greatest common divisor in \(K[X]\).
A nonconstant common factor means that this greatest common divisor has
positive degree; the factor is allowed to depend on \(j\).
Write \(c\ne0\) for the leading coefficient of \(f\).
\paragraph{Statement recorded in the supplied source.}
For every characteristic-zero field \(K\) and every polynomial
\(f\in K[X]\) of degree \(d\ge1\), does
\[
 \deg\gcd(f,f^{(j)})\ge1\quad\text{for every }1\le j\le d-1
\]
imply that there is an \(a\in K\) with \(f(X)=c(X-a)^d\)?
Equivalently, after division by its leading coefficient, must \(f\)
be a \(d\)-th power of a monic linear polynomial? \sref{1107}{2}
\subsection{Short English statement}
If a polynomial shares a root with each of its nonconstant-order derivatives, must all its roots coincide? Different derivatives may share different roots a priori.
\subsection{Sources}
\begin{description}
\item[\hypertarget{source:1107:1}{S1}] ULAM.AI and the UnsolvedMath Contributors, \emph{UnsolvedMath: A Curated Collection of Open Mathematics Problems}, record 1107 (ALG-007). CC BY 4.0. \url{https://huggingface.co/datasets/ulamai/UnsolvedMath}.
\item[\hypertarget{source:1107:2}{S2}] Soham Ghosh, \emph{A finiteness result towards the Casas--Alvero Conjecture}, arXiv:2402.18717, revised version, Section 1, Conjecture CA. \url{https://arxiv.org/abs/2402.18717}.
\end{description}
\label{end:1107}
\end{document}
